\printedpage{64}
% head: A. Grothendieck : Dérivateurs, Ch. I, Édition des Universités de Montpellier II et Paris VII
% foot: Version 8 juillet 2001

\msp{page 107}

The question arises of finding $f$ which are cohomological $\DD$-fibrations, or homological ones. I do not know whether I shall be able to find any, which are not already cohomological \HOT-fibrations, or homological ones (this now amounts to the same thing\footnote[55]{Cohomological \HOT-equivalences $=$ homological \HOT-equivalences $=$ \HOT-equivalences plain and simple.} ), for in practice I hardly know how to say that an arrow $Z' \to Z$ in $\Cat$ is a complete cohomological $\DD$-equivalence, except when this arrow is a \HOT-equivalence (and assuming axiom \textbf{Der} 3 $b'_1$), or its weakened cohomological variant).

I shall therefore look for the \HOT-fibrations, \textit{i.e.} the $f : X \to Y$ such that for every $Y'' \to Y' \to Y$, with $Y''$ and $Y'$ aspheric (understood: for \HOT), $X'' \to X'$ is a \Hot-equivalence. (It even suffices to suppose $Y''$ reduced to a point, $Y'$ with a final object $\ldots$ )

Suppose for example $f$ proper (\textit{i.e.} \HOT-proper), so condition $a$) of definition 2, p. 101 is satisfied and likewise after every base change. Does it suffice to check in addition condition $b$), \textit{i.e.} that
\[ X/y \longrightarrow X/y' \quad , \qquad \text{for } y \to y' \text{ in } Y \ , \]
are \Hot-equivalences? (I mean staying over $Y$, without also supposing we have the same thing after base change).
\msp{page 108}
But does this remain true after every base change? % [sic] reste-il -> reste-t-il, forme interrogative

The situation seems better here for smooth $f$ (instead of proper), interpreting conditions $a$) and $b$) in homological form, as meaning that $f_!(e_X)$ is a \underline{locally} \underline{constant} object of $\mathsf{HOT}(Y)$, which is computed fiber by fiber. Since smooth $f$ implies that $f_!$ commutes with every base change, this hypothesis on $f$ (that $f$ be a cohomological \HOT-fibration, perhaps not complete) is stable under base change, hence implies that it is a complete cohomological \HOT-fibration. But this argument uses the general properties of derivators for \HOT, which I have not yet established - especially not in the whole domain $\Cat$! So it is healthy for me to be wary. I prefer to feel out the difficulty while staying at first within the framework of the \underline{cohomological} formalism (noncommutative), of which I have some mastery. I therefore suppose $f$ proper again, and that the % [sic] meilleur -> meilleure, accord avec « situation »
\[ X_y \longrightarrow X/y \quad , \qquad X/y \longrightarrow X/y' \]
\msp{page 109}
are \HOT-equivalences, I want to prove the same thing after every base change, at least that in the situation
\[
\begin{tikzcd}
X \arrow[r, leftarrow, "q"] \arrow[d, "f"'] & X' \arrow[r, leftarrow, "q'"] \arrow[d, "f'"'] & X'' \arrow[d, "f''"] \\
Y \arrow[r, leftarrow, "p"'] & Y' \arrow[r, leftarrow, "p'"'] & Y''
\end{tikzcd}
\]
\printedpage{65}
% head: A. Grothendieck : Dérivateurs, Ch. I, Édition des Universités de Montpellier II et Paris VII
% foot: Version 8 juillet 2001
if $Y'$, $Y''$ are aspheric, then $q'$ is a \Hot-equivalence. If I want to establish this by a cohomological criterion, I must start by taking on $X'$ a locally constant sheaf of \textit{sets} $F$, and prove that $\Gamma(X',F) \overset{\sim}{\to} \Gamma(X'',{q'}^{*}F)$. Since this sheaf does not come from $X$, I am at a loss how to apply the hypothesis on $f$, which concerns a priori constant sheaves of sets on $X$, or on the $X/y$. But we should have this:

\textbf{Lemma.} \textit{Let $f : X'' \to X'$. Then $f$ is a} \HOT-\textit{equivalence} i.e.
\[ H^{*}_{\mathsf{HOT}}(X', a_{X'}) \xrightarrow{\ \sim\ } H^{*}_{\mathsf{HOT}}(X'', a_{X''}) \]
\textit{is} \texttt{[an]} \textit{isomorphism for every $a \in \mathsf{Hot} = \mathsf{HOT}(e)$, if and only if $f$ is a} \Hot-\textit{equivalence,}\footnote[56]{Note here \texttt{[illegible word]} the distinction between \HOT-equivalence and \Hot-equivalence.} \textit{and then $f$ is even a} $2$-\textit{complete} \HOT-\textit{equivalence} (I need only: complete cohomological and homological \HOT-equivalence) i.e. \textit{for every $\xi \in \DD(X')$ \underline{locally} \underline{constant} on $X'$, we have}\footnote[57]{Is this also true for the derivators $\DD_i$?}
\[ H^{*}_{\mathsf{HOT}}(X',\xi) \xrightarrow{\ \sim\ } H^{*}_{\mathsf{HOT}}(X'',{q'}^{*}(\xi)) \quad . \]
\msp{page 110}

The first assertion should follow from the formulas
\[ H^{*}_{\mathsf{HOT}}(Z,a_Z) \simeq \underline{\mathsf{Hom}}_{\mathsf{Hot}}(\mathsf{hot}(Z),a) \tag*{$(*)$} \]
\[ \text{(and also} \qquad\qquad H_{*}^{\mathsf{HOT}}(Z,a_Z) \simeq \mathsf{hot}(Z) \times a \qquad\qquad ) \]
hence if
\[ \mathsf{hot}(X'') \xrightarrow{\ \sim\ } \mathsf{hot}(X') \tag*{$(**)$} \]
is \texttt{[an]} isomorphism, then
\[ H^{*}_{\mathsf{HOT}}(X', a_{X'}) \longrightarrow H^{*}_{\mathsf{HOT}}(X'', a_{X''}) \]
is \texttt{[an]} isomorphism, since it is deduced from $(*)$ by taking $\underline{\mathsf{Hom}}_{\mathsf{Hot}}(\,.\,,a)$. Conversely, if $(*)$ is always \texttt{[an]} isomorphism, expressing the two sides \texttt{[illegible word]} of $\underline{\mathsf{Hom}}_{\mathsf{Hot}}(\,.\,,a)$ and taking the $\pi_0$, we find that
\[ \mathsf{Hom}_{\mathsf{Hot}}(\mathsf{hot}(X'),a) \xrightarrow{\ \sim\ } \mathsf{Hom}_{\mathsf{Hot}}(\mathsf{hot}(X''),a) \]
is an isomorphism for every $a$, hence $(**)$ is an isomorphism.

But the fact that this implies that $f''$ is a complete cohomological \HOT-equivalence, \textit{i.e.} that the relation $(*)$ remains true for a $\xi \in \DD(X')$ \underline{locally} \underline{constant} on $X'$, seems to me a delicate property of the derivator \HOT, to which I shall have to return.

So the lemma is far from established, for want of having so much as constructed the derivator \HOT, without even speaking of establishing its delicate special properties! But I have no reason to doubt the lemma, it is part of the litany of basic intuitions that I would like to put on a solid footing.

\printedpage{66}
% head: A. Grothendieck : Dérivateurs, Ch. I, Édition des Universités de Montpellier II et Paris VII
% foot: Version 8 juillet 2001

\msp{page 111}

The lemma implies that a $\DD = \mathsf{HOT}$-fibration (in the sense of definition 2 page 81, revised and corrected on page 96-97) is in fact already a complete cohomological and homological \HOT-fibration. Now we have seen that if $f$ is \underline{proper}, $f$ is a $\DD$-fibration if and only if the preceding conditions \underline{over $Y$} (without base change) are satisfied. And dually for smooth $f$.

\begin{flushright}
11.11.90
\end{flushright}

I feel the need to review the notions of $\DD$-fibrations of all kinds (p. 76 ff). First the various useful variants of axiom \textbf{Der} 3 on derivators:

\textbf{Der} 3 $a$)\footnote[58]{\textit{cf.} also the series \textbf{Der 3} $c \ldots$ , page 124 ff.} If $Z \in \mathsf{Dia}$ is \Hot-aspheric, it is $\DD$-aspheric \textit{i.e.} $\mathcal{A} = \DD(e) \to \DD(Z)$ is fully faithful. % Doubt resolved: attached PDF leaf 30 (printed p.66), footnote 58 has bold Der 3.

$a'$) As $a$), with in addition the characterization of the essential image of $\mathcal{A}$ in $\DD(Z)$, as consisting of the locally constant $\DD$-coefficients. % [sic] localements -> localement, adverbe invariable

\textbf{Der} 3 $b$) If $f : Z' \to Z$ in \Dia{} is a \Hot-equivalence, then it is a $\DD$-equivalence, \textit{i.e.} the functor $f^{*} : \DD(Z) \to \DD(Z')$ induces an equivalence between the subcategories consisting of the constant coefficients (or again a fully faithful functor $\DD(Z)_{\text{loc. const.}} \longrightarrow \DD(Z')$).

$b'_1$) If $f$ is as above, then $f$ is a 1-complete $\DD$-equivalence, \textit{i.e.} the functor $f^{*} : \DD(Z) \to \DD(Z')$ induces a fully faithful functor
\msp{page 112}
between the full subcategories consisting of the locally constant objects.

$b'_2$) As $b'_1$), but with 2-\underline{complete} $\DD$-equivalence instead of 1-complete, \textit{i.e.} $f^{*}$ induces an equivalence between the subcategories consisting of the locally constant coefficients.

$b''_{\mathrm{coh}}$) If $f : Z' \to Z$ in \Dia{} is a \Hot-equivalence, then it is a complete cohomological $\DD$-equivalence, \textit{i.e.} for every $F \in \mathsf{Ob}\ \DD(Z)$ which is \underline{locally} \underline{constant}, we have
\[ H^{*}_{\DD}(Z,F) \xrightarrow{\ \sim\ } H^{*}_{\DD}(Z',f^{*}F) \quad . \]

$b''_{\mathrm{hom}}$) If $f : Z' \to Z$ in \Dia{} is a \Hot-equivalence, it is a complete homological $\DD$-equivalence, \textit{i.e.} for every $F \in \mathsf{Ob}\ \DD(Z)$ which is \underline{locally} \underline{constant}, we have\footnote[59]{\textit{cf.} special $\pm$ variant for \HOT{} pages 124-125.}
\[ H_{*}^{\DD}(Z',f^{*}F) \xrightarrow{\ \sim\ } H_{*}^{\DD}(Z,F) \quad . \]

\printedpage{67}
% head: A. Grothendieck : Dérivateurs, Ch. I, Édition des Universités de Montpellier II et Paris VII
% foot: Version 8 juillet 2001

We have the diagram of tautological implications for the variants of \textbf{Der} 3:
\[
\begin{tikzcd}
 & b'_2 \arrow[d, Rightarrow] \arrow[drr, Rightarrow] & & & \\
 & b'_1 \arrow[dl, Rightarrow] \arrow[dr, Rightarrow] \arrow[dd, Rightarrow] & & a' \arrow[ddd, Rightarrow] & \\
b''_{\mathrm{coh}} \arrow[dr, Rightarrow] & & b''_{\mathrm{hom}} \arrow[dl, Rightarrow] & & \\
 & b \arrow[drr, Rightarrow] & & & \\
 & & & a & .
\end{tikzcd}
\]
In all cases known to me, derivators satisfy the strongest condition \textbf{Der} 3 $b'_2$.
\msp{page 113}

I also recall the:

\textbf{Definition 1:} Let $f : X \to Y$ be in \Dia. We say that $f$ is a \textit{cohomological $\DD$-equivalence} (understood: ``\textit{complete}'') if for every $\DD$-coefficient $F$ \underline{locally} \underline{constant} on $Y$,
\[ H^{*}_{\DD}(Y,F) \xrightarrow{\ \sim\ } H^{*}_{\DD}(X,f^{*}F) \quad ; \]
we say that $f$ is a \textit{homological $\DD$-equivalence} (``\textit{complete}'') if for every $F$ as \texttt{[here-]}above,
\[ H_{*}^{\DD}(X,f^{*}F) \xrightarrow{\ \sim\ } H_{*}^{\DD}(Y,F) \quad . \]
We say that $f$ is a 1-\textit{complete} $\DD$-\textit{equivalence} (resp. 2-\textit{complete}) if the functor
\[ f^{*} : \DD(Y)_{\text{loc. const.}} \longrightarrow \DD(X)_{\text{loc. const.}} \]
is fully faithful (resp. an equivalence).
\[
\begin{tikzcd}[column sep=small]
 & \text{2-complete $\DD$-equivalence} \arrow[d, Rightarrow] & & \\
 & \text{1-complete $\DD$-equivalence} \arrow[dl, Rightarrow] \arrow[dr, Rightarrow] & & \\
\text{cohomological $\DD$-equivalence} \arrow[dr, Rightarrow] & & \text{homological $\DD$-equivalence} \arrow[dl, Rightarrow] & \\
 & \text{$\DD$-equivalence} & & .
\end{tikzcd}
\]
If in the subcategories $\DD(X)_{\text{loc. const.}}$ there is an operation $*$ (associative, unital, commutative) and associated internal \underline{\textsf{Hom}}, preserved by inverse images, then cohomological $\DD$-equivalences are already 1-complete and are therefore homological $\DD$-equivalences, therefore in this particular case
\printedpage{68}
% head: A. Grothendieck : Dérivateurs, Ch. I, Édition des Universités de Montpellier II et Paris VII
% foot: Version 8 juillet 2001
\msp{page 114 vide}
\msp{page 115}
\[
\begin{tikzcd}
\text{2-complete $\DD$-equivalence} \arrow[d, Rightarrow] & & \\
\text{1-complete $\DD$-equivalence} \arrow[r, Leftrightarrow] & \text{cohomological $\DD$-equivalence} \arrow[d, Rightarrow] & \\
 & \text{homological $\DD$-equivalence} \arrow[d, Rightarrow] & \\
 & \text{$\DD$-equivalence} & .
\end{tikzcd}
\]
In $\DD = \mathsf{HOT}$, all these notions should be equivalent, and equivalent to the apparently still weaker condition, that $f : X \to Y$ be a \Hot-equivalence. But I know only \HOT{} as a derivator where the $\DD$-equivalence of an arrow $f : X \to Y$ in \Dia{} already implies cohomological $\DD$-equivalence (complete), or homological $\DD$-equivalence (complete), or even 1-complete $\DD$-equivalence. (I shall come back to this presently.)

\textbf{Proposition 1.} \textit{Let $f : X \to Y$ be in} \Dia. \textit{Equivalent conditions:}

a) \textit{For every $y \in \mathsf{Ob}\, Y$, the inclusion $X_y \to X/y$ is a $\DD$-equivalence, and for every arrow $y' \to y$ in $Y$, $X/y' \to X/y$ is a $\DD$-equivalence.}

b) \textit{For every \underline{constant} coefficient $F$ on $Y$, $f_{*}f^{*}F$ is locally constant, and is computed ``fiber by fiber''.}

(This is prop. 1, p. 76).

\textbf{Proposition 1$'$.} Dual.

\msp{page 116}

\textbf{Definition 2.} Under the equivalent conditions of proposition 1, we say that $f$ is a \textit{weak cohomological $\DD$-fibration.} Under the dual conditions, that $f$ is a \textit{weak homological $\DD$-fibration.}

This is what in definition 1, p. 77 I called ``cohomological $\DD$-fibration'' or ``left $\DD$-fibration'', and dually. Now the notion seems to me too weak to deserve this name, which I reserve for the ``complete'' case.

\textbf{Proposition 2.} \textit{Weak cohomological $\DD$-fibrations are stable under smooth base changes, weak homological $\DD$-fibrations stable under proper base changes.}

I should construct an example where $f$ is a weak cohomological $\DD$-fibration, and is no longer so after base change.